\documentclass[10pt,a4paper]{amsart}
\usepackage[margin=19mm]{geometry}
\usepackage{amsmath,amssymb,mathtools}
\usepackage[scr=rsfs,cal=euler]{mathalfa}
\usepackage{xfrac}
\usepackage[unicode,hidelinks,bookmarks=false]{hyperref}
\newcommand{\ie}{i.e.\ }
\newcommand{\cf}{cf.\ }
\newcommand{\resp}{respectively\ }
\newcommand{\Subsection}{\subsection}
\providecommand{\nobreakdash}{\nobreak\hbox{-}}
\setlength{\emergencystretch}{2em}
\multlinegap0pt
\usepackage{bm}
\usepackage{bbm}
\usepackage{dsfont}
\usepackage{xspace}
\usepackage{cancel}
\usepackage{mathtools}
\usepackage{amsthm}
\makeatletter
\@ifundefined{theorem}{\newtheorem{theorem}{Theorem}}{}
\makeatother
\newcommand{\lto}{\longrightarrow}
\newcommand{\sD}{\mathscr{D}}
\newcommand{\sR}{\mathscr{R}}
\newcommand{\sS}{\mathscr{S}}
\newcommand{\sT}{\mathscr{T}}
\newcommand{\sU}{\mathscr{U}}
\title[Continuity of the superpotentials and slices of tropical cur]{Continuity of the superpotentials and slices of tropical currents}
\author{Farhad Babaee, Tien Cuong Dinh}
\date{Source: arXiv:2506.09828v1 (CC BY 4.0)}
\begin{document}
\maketitle

\noindent\emph{Source and licence.} Continuity of the superpotentials and slices of tropical currents. Version arXiv:2506.09828v1 (CC BY 4.0), distributed under the Creative Commons Attribution 4.0 International licence, as recorded in the arXiv licence metadata for this exact version and shown on the abstract page https://arxiv.org/abs/2506.09828v1. Full rights notice: licenses/arxiv-2506.09828-CC-BY-4.0.txt.\par\smallskip

\noindent\textit{Editorial scope.} Counted unit: the Theorem whose source label is \texttt{thm:CS-blowup} in the article (source0.tex lines 490--492, source label \texttt{thm:CS-blowup}), delivered together with the article's own complete original proof of it (source0.tex lines 493--528, one \texttt{proof} environment of 36 lines). No mathematics is written, repaired or re-derived: statement and proof are the article's own text line for line. Required context retained uncounted: none. Every definition, display and label of those lines is retained unchanged. Omitted from this package and not redistributed: 1--489, 529--2524. Nothing inside a retained excerpt is omitted or altered: every retained body is delivered exactly as the article's lines are written, so no omission receipt of any kind accompanies this package. Citations: the retained excerpts carry no citation command at all, so no bibliography entry of the article is transcribed and no reference is redistributed. Reference adaptation: none was needed and none was made; every reference inside the delivered unit resolves inside it, so no unresolved reference or citation is produced. Printed slips kept literal: no printed word, symbol, hypothesis or number of the counted statement or of its proof was altered. Layout and class: the article's own class is \texttt{amsart} with packages including bm, amsmath, amsthm, amssymb, mathrsfs, epic, eepic, paralist, enumerate, xy, tikz-cd, xcolor, graphicx, euscript; the wrapper is the lane's pinned \texttt{amsart} editorial wrapper at 10pt on A4, which restates the theorem-like environments the retained text needs and adds the article's own macro definitions that the retained text needs (\texttt{\textbackslash lto}, \texttt{\textbackslash sD}, \texttt{\textbackslash sR}, \texttt{\textbackslash sS}, \texttt{\textbackslash sT}, \texttt{\textbackslash sU}), with extra wrapper packages (bm, bbm, dsfont, xspace, cancel, mathtools). The delivered numbering is the unit's own. Assets: no retained span contains a figure environment, a graphics-inclusion command, a PDF-page inclusion, a table or a TikZ picture; no image file is copied into this package, the package declares no asset dependency and no image file is redistributed with it. Licence: version arXiv:2506.09828v1 of this article is distributed under the Creative Commons Attribution 4.0 International licence, as recorded in the arXiv licence metadata for this exact version and shown on the abstract page https://arxiv.org/abs/2506.09828v1. A full rights notice is delivered at licenses/arxiv-2506.09828-CC-BY-4.0.txt. Characters: every retained excerpt is pure ASCII, so the delivered engine, XeLaTeX with its default Latin Modern text font, reports no missing character.
\begin{theorem}\label{thm:CS-blowup}
  Let $q: \widehat{X} \lto X,$ be the blowing up of the compact K\"ahler manifold $X$ along a submanifold. Assume that $\sT \in \sD_p(\widehat{X})$ is such that the support of $\sT$ does not intersect the exceptional divisors of $\widehat{X}.$ If   the   current $\sT$ has a continuous superpotential then $q_* \sT$ has the same property. 
\end{theorem}

\begin{proof} Assume that $\sS\in \sD^{d-p+1}(X)$. We need to show that the superpotential of $q_* \sT$, which is a function defined on smooth forms in $\sD^{d-q+1,0}(X)$, can be extended to a continuous function on 
$\sD^{d-q+1,0}(X)$. Let $\alpha$ be a smooth closed $(p,p)$-form cohomologous to $q_*\sT$ and define $\beta:=q^*\alpha$. By hypothesis, we have $q^*q_*\sT=\sT$. It follows that $\sT$ is cohomologous to $\beta$. Fix a potential $U$ of $q_* \sT-\alpha$ which is smooth near the blowup locus. We have $dd^c  U= q_* \sT-\alpha$ and therefore 
$dd^c q^*U = \sT-\beta$. The smoothness of $U$ near the blowup locus implies that $q_*(q^*U)=U$.

Since $\alpha$ is smooth, it is enough to show that the superpotential $\sU$ of $q_* \sT-\alpha$ can be extended to a continuous function on $\sD^{d-q+1,0}(X)$. Since the last current has a vanishing cohomology class, its superpotential doesn't depend on the normalization. 


\bigskip\noindent
{\bf Claim 1.}  Let $(\sS_n)$ be a bounded sequence of smooth forms in $\sD^{d-p+1}(X)$. Then the sequence $\sU(\sS_n)$ is bounded. 

\proof[Proof of Claim 1]
Since $\sS_n$ is smooth, we have
$$\sU(\sS_n)=\langle U,\sS_n\rangle = \langle q^*(U), q^*\sS_n \rangle = \widehat\sU(q^*\sS_n), $$
where $\widehat \sU$ denotes the superpotential of $\sT-\beta$.
Since the action of $q^*$ on cohomology is bounded and the mass of a positive closed current only depends on its cohomology class, we see that $q^*\sS_n$ is bounded in $\sD^{d-p+1,0}(\widehat X)$. Since $\sT$ has a continuous superpotential, we deduce that $\widehat\sU(q^*\sS_n)$ is bounded and hence $\sU(\sS_n)$ is bounded as claimed.
\endproof


\bigskip\noindent
{\bf Claim 2.}  Let $(\sS_n)$ be a bounded sequence of smooth forms in $\sD^{d-p+1}(X)$ converging to 0. Then $\sU(\sS_n)$ tends to 0.

\proof[Proof of Claim 2]
By extracting a subsequence, we can assume that $q^*\sS_n$ converges to some current $\sR$ supported by the exceptional divisor. 
By hypothesis, we have $q_*(\sR)=0$.
Using the computation in Claim 1, and the fact that $\widehat \sU$ is continuous, we get
$$\lim  \sU(\sS_n) = \widehat\sU (\sR).$$
Now, since $q^*U$ is smooth near the exceptional divisor which supports $\sR$, we deduce that 
$$\widehat\sU (\sR)=\langle q^*U,\sR\rangle = \langle U, q_*\sR\rangle =0.$$
This proves the claim.
\endproof

To finish the proof, let $\sS$ be any current in $\sD^{d-p+1,0}(X)$. Choose a bounded sequence $\sR_n$ of smooth forms in $\sD^{d-p+1,0}(X)$ converging to $\sS$. 
By Claim 1, extracting a subsequence allows to assume that $\sU(\sR_n)$ converges to some real number $l$. Consider now an arbitrary bounded sequence of 
$(\sS_n)$ of smooth forms in $\sD^{d-p+1,0}(X)$ converging to $\sS$. To show that $\sU$ extends to a continuous function at $\sS$, it is enough to check 
that $\sU(\sS_n)$ converges to $l$. This is a consequence of Claim 2 applied to the sequence $\sS_n-\sR_n$.
\end{proof}

\end{document}
